\section{The Archi Agent}
\label{sec:agent}

Archi is implemented as an open-source Python framework
configured per deployment; this section describes the CompOps
configuration. Section~\ref{sec:components} covers the three
top-level components, Section~\ref{sec:pipelines} the agent
pipeline and the controls evaluated against it,
Section~\ref{sec:tools} the tool catalogue, and
Section~\ref{sec:deployment} the production rollout.

\subsection{Components}
\label{sec:components}

The framework decomposes into three components: a data manager
that ingests and indexes documentation sources, a pipeline layer
that the chat service routes user questions through, and a chat
interface. Figure~\ref{fig:architecture} shows the top-level
dataflow.

\begin{figure}[htbp]
\centering
\begin{tikzpicture}[
    >={Latex[length=2.2mm]},
    every node/.style={font=\small},
    node distance=7mm,
    box/.style={
        rectangle, rounded corners=2pt, draw, line width=0.7pt,
        inner sep=3pt, align=center, minimum height=9mm,
    },
    iobox/.style={box, fill=gray!8, minimum width=58mm},
    rxn/.style={box, fill=agentfill, draw=agentblue, line width=1.0pt,
        minimum width=42mm},
    docbox/.style={box, fill=docfill, draw=docgreen, line width=0.7pt,
        minimum width=27mm, minimum height=13mm},
    livebox/.style={box, fill=livefill, draw=liveorange, line width=0.7pt,
        minimum width=27mm, minimum height=13mm},
    arrlbl/.style={font=\footnotesize\itshape, inner sep=1.5pt,
        fill=white, inner ysep=0.5pt}
]

% Top: question
\node[iobox] (q) {Operator question};

% ReAct loop core
\node[rxn, below=of q] (reason)
    {\textcolor{agentblue}{\textbf{Reason}} (LLM step)};
\node[rxn, below=of reason] (tool)
    {\textcolor{agentblue}{\textbf{Tool call}}};

% Tool families
\node[docbox, below=of tool, xshift=-15mm] (docs)
    {\textcolor{docgreen}{\textbf{Doc.\ retrieval}}\\
     \footnotesize TWiki, Jira};
\node[livebox, below=of tool, xshift=15mm] (live)
    {\textcolor{liveorange}{\textbf{Live MonIT}}\\
     \footnotesize Rucio, HTCondor};

% Final answer below, centered under the loop
\node[iobox, below=24mm of tool] (out)
    {Answer with tool-call trace};

% --- Forward flow ---
\draw[->, line width=0.8pt] (q.south) -- (reason.north);
\draw[->, line width=0.8pt] (reason.south) -- (tool.north)
    node[arrlbl, midway, right=1pt] {dispatch};
\draw[->, line width=0.8pt] (tool.south) -- (docs.north);
\draw[->, line width=0.8pt] (tool.south) -- (live.north);

% --- Loop back: tool return on the right ---
\draw[->, line width=0.8pt, rounded corners=4pt]
    (tool.east) -- ++(11mm,0) |- (reason.east)
    node[arrlbl, pos=0.25, right=1pt] {return};

% --- Final exit: reason directly emits the answer (curves down on the left) ---
\draw[->, line width=0.8pt, rounded corners=4pt]
    (reason.west) -- ++(-11mm,0) |- (out.west)
    node[arrlbl, pos=0.04, left=1pt] {final};

\end{tikzpicture}
\caption{Archi's ReAct loop. Each iteration the
\textsc{Reason} step dispatches a \textsc{Tool call} (either
documentation retrieval or live MonIT), and the tool return
becomes the next reason input. When \textsc{Reason} decides no
further tools are needed it emits the answer with its tool-call
trace.}
\label{fig:architecture}
\end{figure}

The data manager ingests web pages, git repositories, local
files, and Jira projects into a PostgreSQL \texttt{pgvector}
index of chunked passages embedded with
\texttt{sentence-transformers/all-MiniLM-L6-v2}~\cite{reimers2019sbert}.
For the CompOps deployment the sources are the CMS TWiki, a
curated list of public CMS documentation repositories, and seven
CompOps Jira projects (\texttt{CMSPROD}, \texttt{CMSTRANSF},
\texttt{CMSDM}, \texttt{CMSTZ}, \texttt{PRCAMPAIGNS},
\texttt{CMSMONIT}, \texttt{CMSVOC}); ticket text is anonymised at
ingest.

A pipeline takes a question, the prior conversation, and the
deployment configuration, and returns an answer with the
documents and tool calls that supported it. Pipelines are
declared in YAML and selected per deployment; the pipelines
compared in Section~\ref{sec:evaluation} differ on two axes:
whether they retrieve from the documentation store and whether
they call live MonIT.

\subsection{The CompOps Agent}
\label{sec:pipelines}

The deployed assistant runs an agent built on the ReAct pattern
\cite{yao2023react} using the LangGraph state-graph framework
\cite{langgraph}.
The graph alternates a reasoning step (which emits a textual
trace) with a tool-call step (which invokes either a
documentation-retrieval tool or a live MonIT query); when the
reasoning step decides no further tools are needed, it emits the
final answer with the accumulated tool-call trace.
The system prompt steers the agent toward parallel exploratory
tool calls, aggregation queries when counting, and fetch-by-ID
fallback when a search returns only summaries. It also forbids
asking the operator for clarification, requiring a best-guess
answer with explicit refusal when live state is unreachable.

Two control pipelines and one ablation are summarised in
Table~\ref{tab:pipeline_summary}. The bare-LLM control sends the
question and conversation history directly to the language model;
the RAG control retrieves the top-eight passages from the
documentation store using a hybrid BM25-plus-vector retriever and
conditions on them~\cite{lewis2020rag}; the \emph{agent (no live
tools)} ablation runs the agent with the live MonIT tools removed
while retaining documentation retrieval, isolating the
contribution of live MonIT relative to the full agent and the
contribution of the agent loop relative to RAG.

\begin{table}[htbp]
\centering
\caption{The four pipelines compared in this paper, including a
no-live-tools ablation of the agent. Documentation refers to the
hybrid retriever over the CMS TWiki and ingested Jira tickets;
MonIT refers to live OpenSearch queries against the Rucio and
HTCondor monitoring indices.}
\label{tab:pipeline_summary}
\begin{tabular}{lcc}
\toprule
\textbf{Pipeline} & \textbf{Documentation} & \textbf{Live MonIT} \\
\midrule
Bare LLM               & no  & no  \\
RAG                    & yes & no  \\
Agent (no live tools)  & yes & no  \\
Agent                  & yes & yes \\
\bottomrule
\end{tabular}
\end{table}

\subsection{Tools}
\label{sec:tools}

The agent has access to three tool families, listed in
Table~\ref{tab:tool_catalogue}: documentation retrieval, the
Rucio data-management MonIT index, and the
HTCondor~\cite{htcondor2005} jobs MonIT index. Each MonIT family
exposes search, aggregation, and fetch variants; the documentation
family exposes hybrid BM25-plus-vector search, metadata-index
browse, and fetch-by-ID. MonIT tools authenticate against the
central CERN Grafana-OpenSearch endpoint with bearer tokens; every
tool call is recorded in the trace surfaced to the operator.

\begin{table}[htbp]
\centering
\caption{Tools exposed to the agent. Each family supports a
search, an aggregation, and a fetch-by-ID variant against its
backend.}
\label{tab:tool_catalogue}
\begin{tabular}{@{}lp{0.55\columnwidth}@{}}
\toprule
\textbf{Family} & \textbf{Backend} \\
\midrule
Documentation  & \texttt{pgvector} over CMS TWiki and Jira tickets \\
Rucio MonIT    & WLCG OpenSearch, Rucio events \\
HTCondor MonIT & WLCG OpenSearch, HTCondor jobs \\
\bottomrule
\end{tabular}
\end{table}

\subsection{Production Deployment}
\label{sec:deployment}

Archi has run as the CMS CompOps operator chat assistant since
February 2026 on shared CMS infrastructure at CERN, serving the
operations team, on-call experts, and contributors debugging
production workflows. Operators interact through a web chat interface that
maintains turn-by-turn conversation history, lists source
documents per response, and exposes the agent's tool-call trace.
The deployed pipeline runs on GPT-5; the evaluation in
Section~\ref{sec:evaluation} re-scores those production answers
on the same rubric and benchmarks local-weight alternatives
against them.
